\section{Future work}
\label{sec:future}

Further work should build on the locked-window, mission, and ablation findings rather than repeating those evaluations.

\subsection{Radiation across activity regimes}

Radiation is the highest priority because the locked hybrid did not outperform the replicated public baseline in either window. Future studies should cover quieter and stormier regimes, denser OMNI/DONKI/SWPC sampling, improved proton-flux treatment, calibrated uncertainty for alerts, and transparent comparison with combined NASA warning paths without retrospective threshold tuning.

\subsection{Solar input behaviour}

The solar ablation result should be repeated under stronger activity and, where feasible, with retrained rather than mean-substituted configurations, to show when live solar input helps most.

\subsection{Independent labels for physics-tied channels}

Temperature, moonquakes, dust, and micrometeorites should remain physics-led, but scoring needs targets independent of the physics path. Micrometeorites in particular should be assessed against observations that do not already contain the LIPAS shower term.

\subsection{Validation with measured data}

As Artemis and other surface or orbital measurements become available, evaluation labels should be replaced with direct measurements wherever possible. Collaboration with space agencies and research groups can strengthen validation criteria.

\subsection{Operational development}

Interface and mission-planner work should improve uncertainty visualisation, checklist-aligned reports, multi-hazard simulations, and tests of behaviour when data feeds fail. Usability should eventually be measured with mission personnel.

\subsection{Beyond the lunar surface}

Because much of the system already uses heliophysics inputs, extension to cislunar space and Earth orbit is a natural next step, with Mars as a longer-term possibility.

Throughout development, LIPAS should remain consultative decision support used alongside NASA and NOAA products, not a replacement for specialised or flight-qualified systems.
